% figures.tex -- Figures collected after the tables
\section*{Figures}

\begin{center}
\begin{minipage}{\linewidth}
\captionsetup{type=figure,hypcap=false}
  \centering
  \includegraphics[width=0.95\linewidth]{figs/Fig1-RuleQuench.pdf}
  \caption{The change in traders' decision rules and its effect on the same starting
market. \textbf{(A)}~Original Rule~19280. \textbf{(B)}~New Rule~19286. Each
cross contains the four neighbors whose current decisions determine the
trader's next decision, shown below the arrow. The hollow center marks the
trader being updated; its current decision does not enter the rule. Blue,
gray, and orange denote Buy~(0), Hold~(1), and Sell~(2). One neighbor
arrangement is shown for each possible sum; all arrangements with that sum
give the same response. The outlined column is the only change: three
neighbors in Buy and one in Hold lead to Buy under Rule~19280 and Sell under
Rule~19286. \textbf{(C)}~Decisions on a $64\times64$ grid when the new rule
was assigned and after 100, 400, and 3,000 steps. The upper row keeps the
original rule throughout the interior; the lower assigns the new rule to
20\% of interior traders. Both markets share the same 1,000-step history,
unlimited dealer, and active technical traders. Seed 12 was
selected from 50 starting markets by proximity to the median difference in
peak Sell fraction between changed and unchanged markets. Outcomes across
starting markets are reported separately (Supplementary
Table~\ref{tab:paired-switch}).}
  \label{fig:rule-quench}
\end{minipage}
\end{center}

\clearpage

\begin{center}
\begin{minipage}{\linewidth}
\captionsetup{type=figure,hypcap=false}
  \centering
  \includegraphics[width=0.97\linewidth]{figs/Fig2-ReservoirQuench.pdf}
  \caption{Changing decision rules while varying how much the dealer could buy and
sell. \textbf{(A)}~Probability of a Sell majority with the technical traders kept
in Hold (held) and with them trading on the price (active) at selected dealer
capacities. Bands are 95\% Wilson intervals over 50 starting
markets. \textbf{(B)}~Difference in majority probability between the active
and held conditions, across all capacities. \textbf{(C)}~Median maximum drawdown at
$f=0.20$; bands mark the interquartile range. \textbf{(D)}~Median fraction
of orders backed by trader cash or shares that the dealer rejected at
$f=0.20$. Capacity $h=10$ appears as a dashed overlay in (A) and a separate
capacity in (B--D); it matched $h=\infty$ exactly.}
  \label{fig:reservoir}
\end{minipage}
\end{center}

\clearpage

\begin{center}
\begin{minipage}{\linewidth}
\captionsetup{type=figure,hypcap=false}
  \centering
  \includegraphics[width=0.97\linewidth]{figs/Fig3-FiniteSizeQuench.pdf}
  \caption{Effect of market size on the formation of a Sell majority over 50 starting
markets. \textbf{(A)}~Majority probability with the technical traders kept in Hold;
bands are 95\% Wilson intervals. \textbf{(B)}~Variation between runs with
the technical traders kept in Hold, measured as susceptibility: the variance of the
final 500-step mean Sell fraction multiplied by the number of interior
traders. \textbf{(C)}~Fraction $f_{50}(L)$ at which half the runs reached a
majority, estimated by interpolation with the technical traders kept in Hold
and for two dealer capacities with them active.
\textbf{(D)}~Maximum susceptibility versus grid size on linear axes. Error
bars in (C,D) are 95\% intervals from resampling starting markets. Solid
lines connect tested values; dashed curves in (A) are logistic fits. No law
for the limit of indefinitely large grids was fitted.}
  \label{fig:finite-size}
\end{minipage}
\end{center}

\clearpage

\begin{center}
\begin{minipage}{\linewidth}
\captionsetup{type=figure,hypcap=false}
  \centering
  \includegraphics[width=0.90\linewidth]{figs/Fig4-DoseMarketPanels.pdf}
  \caption{How traders' decisions changed over time when different fractions received
 the new rule, with the technical traders active and an unlimited dealer.
Blue, gray, and orange cells denote Buy, Hold, and Sell.
\textbf{(A)}~Decisions in the same five markets (columns) when the rule
changed and after 1,500, 3,000, 5,000, and 10,000 steps (rows). Selection
used the original 3,000-step results and preceded inspection of images.
At that endpoint, the first two markets had no Sell majority and the other
three did. Runs with a majority were selected by proximity to their group's
median time to a majority; runs without one were selected by proximity to
their group's median final Sell fraction. Seeds were 8, 5, 22, 46, and 21
from left to right. The second market subsequently reached a majority at
4,117 steps. \textbf{(B)}~Price divided by its value when the rule changed.
\textbf{(C)}~Interior Sell fraction; the horizontal line marks one-half.
The dark vertical dashed line in (B,C) marks the original 3,000-step endpoint;
other vertical guides mark the remaining snapshots.}
  \label{fig:dose-panels}
\end{minipage}
\end{center}

\clearpage

\begin{center}
\begin{minipage}{\linewidth}
\captionsetup{type=figure,hypcap=false}
  \centering
  \includegraphics[width=0.92\linewidth]{figs/Fig5-DealerMarketPanels.pdf}
  \caption{How the same starting market developed with different dealer capacities at
$f=0.20$. \textbf{(A)}~Decisions under four conditions (columns) when the
rule changed and after 500, 1,500, and 3,000 steps (rows). The first three
columns had the technical traders active, with dealer capacities
$h=0.1$, $h=1$, and $h=\infty$. The last kept the technical traders in Hold
with $h=\infty$. Seed 25 was selected from runs that reached a majority at all
three capacities by minimizing the sum of distances from each capacity's
median time to a majority. Rows show the same elapsed times.
\textbf{(B)}~Price divided by its value when the rule changed.
\textbf{(C)}~Interior Sell fraction. Keeping the technical traders in Hold
prevented price from influencing decisions inside the grid; trades and price changes
were still calculated. This example follows one starting market, while the
full comparison measures outcomes across starting markets
(Fig.~\ref{fig:reservoir}).}
  \label{fig:dealer-panels}
\end{minipage}
\end{center}

\clearpage

\begin{center}
\begin{minipage}{\linewidth}
\captionsetup{type=figure,hypcap=false}
  \centering
  \includegraphics[width=0.99\linewidth]{figs/Fig6-BalanceSheetQuench.pdf}
  \caption{Effects of changing traders' cash and shares when the new decision rule was
introduced, with the technical traders active, $f=0.18$, and 50 starting
markets. In A to C, the nine combinations are grouped by cash multiplier,
with share multiplier on the axis and a separate series for each dealer
capacity. \textbf{(A)}~Probability of a Sell majority. The dashed line and
band give the probability and its 95\% Wilson interval with the technical
traders kept in Hold, which were identical in all nine combinations because resource
changes could not affect decisions. \textbf{(B)}~Median maximum drawdown on
a logarithmic axis. Medians of exactly zero appear at the axis floor with
open markers. \textbf{(C)}~Median fraction of orders backed by trader cash
or shares that the dealer rejected. \textbf{(D)}~Majority probability
against median maximum drawdown for all 27 combinations of cash, shares, and
dealer capacity with the technical traders active. The axis and zero
markers follow B, and the reference probability with the technical traders
kept in Hold is repeated.}
  \label{fig:balance-sheet}
\end{minipage}
\end{center}
